% Title, one-page summary, reading guide and table of contents.
\thispagestyle{plain}
\begin{center}
{\color{accentdark}\rule{\linewidth}{0.8pt}}\\[10pt]
{\LARGE\bfseries Totem-implied repo}\\[6pt]
{\large Controlling the front-office repo and detecting bad curves}\\[8pt]
{\normalsize\color{ink2} Methodology note \quad\textbullet\quad Independent price verification, equity
derivatives}\\[2pt]
{\normalsize\color{ink2} October 2026 \quad\textbullet\quad Draft for discussion}\\[4pt]
{\color{accentdark}\rule{\linewidth}{0.8pt}}
\end{center}

\begin{bref}
\textbf{The problem.} Independent price verification implies the repo curve from Totem: the forwards come from
consensus call and put prices, the dividends from a separate Totem service. The two services do not have the same
contributors, and their number changes with maturity; in our example, \vRunNoptShort{} option contributors and
\vRunNdivShort{} dividend contributors at the short end, \vRunNoptLong{} and \vRunNdivLong{} at five years. The implied
repo curve inherits the noise of both panels: it sometimes looks irregular, or jumps from one maturity to the next for
no market reason. Yet this is the curve that, through the sensitivities, controls the front-office (FO) repo positions.

\medskip
\textbf{The three goals of the note.}
\begin{description}[leftmargin=1.1cm,labelwidth=0.9cm,labelsep=0.2cm,itemsep=1pt]
\item[G1] \emph{Understand}: where do these shapes come from, and which ones are just normal noise?
\item[G2] \emph{Control}: ``I have the forward control, I have the dividend control: how much is the repo control?''
Can the repo position be controlled without depending on the dividend contributors?
\item[G3] \emph{Detect}: how can a bad Totem curve be flagged automatically, its cause named, and the flag kept apart
from an FO mismark?
\end{description}

\medskip
\textbf{The answers.}
\begin{description}[leftmargin=1.1cm,labelwidth=0.9cm,labelsep=0.2cm,itemsep=2pt]
\item[A1] The implied repo equals the average repo of the option panel plus a \emph{panel gap}: the difference
between two dividend averages taken over different contributors. For a European stock, this gap is commonly
\vGapOneY\,bp at one year and \vGapFiveY\,bp at five years (\cref{sec-diagnostic}). No method that only sees Totem
data can remove it; one can, however, compute its size and recognise, by their shape, the mechanisms that cause real
anomalies.
\item[A2] Repo line $=$ forward line $-$ dividend line $+$ non-forward line. For a book whose value depends on the
forwards only (European options, forwards, futures), the non-forward line is zero: the total control does not depend on
the dividend contributors, and it is exactly independent of them under full revaluation (\cref{sec-control}). An
exposure vector computed from the FO sensitivities alone shows where the split between repo and dividends matters.
\item[A3] The structural detector judges the Totem--FO difference ``in carry'', against a noise band computed from the
dispersions and contributor counts that Totem publishes, and tests one typical shape per mechanism
(\cref{sec-detection}). Its false alarm rate is fixed (\vFPR{} of clean names); every alarm comes with an explanation in
natural units and a repaired curve. On a synthetic benchmark of \vNTest{} names (\cref{sec-benchmark}), it detects on
average as much as the best alternative tested, B-spline plus Gaussian mixture (\vMeanStruct{} against
\vMeanGmmDiff{}), finds forward outliers far better (\vFwdStruct{} against \vFwdGmmDiff{}), and does not mistake an FO
mismark for a bad Totem curve (\vFoStruct{} alarms against \vFoGmmDiff{}).
\end{description}
\end{bref}

\newpage
\section*{How to read this note}
Each section opens with its goal and closes with a ``Key points'' box. One synthetic name, the \emph{running example},
illustrates the note from start to finish (green boxes): a European stock seen first in a clean market, then under eight
different anomalies. Short proofs follow the statements; the longer ones are in \cref{app-proofs}.

\begin{center}\small
\begin{tabular}{>{\raggedright\arraybackslash}p{5.1cm}c>{\raggedright\arraybackslash}p{8.0cm}}
\toprule
Question & Section & What you will find\\
\midrule
What are we after, and why is it hard? & \ref{sec-problem} & The control chain, the $1/T$ magnifier, the running
example\\
Which notation? & \ref{sec-setup} & Forward, carry difference, Totem data, from Totem to the implied repo\\
\textbf{G1} Where do bad curves come from? & \ref{sec-diagnostic} & The panel gap and its size, contributor
drop-offs, six mechanisms and their signatures, what the data cannot tell\\
\textbf{G2} How to control the repo? & \ref{sec-control} & The control identity, three books, how much the
repo/dividend split matters\\
\textbf{G3} How to detect? & \ref{sec-detection} & The noise model, the typical shapes, the test, the calibration,
the diagnosis of the running example\\
Does it work? And B-splines, mixtures, neural networks? & \ref{sec-benchmark} & The benchmark and the comparison with
the starting ideas\\
What to do in practice? & \ref{sec-reco} & Recommendations, data needed, monthly procedure\\
\bottomrule
\end{tabular}
\end{center}

{\small\setcounter{tocdepth}{2}\tableofcontents}
\newpage
